%
% 4-fundamental.tex
%
% (c) 2025 Prof Dr Andreas Müller
%
\section{Der Fundamentalsatz der Algebra
\label{buch:komplex:section:fundamental}}
\kopfrechts{Der Fundamentalsatz der Algebra}

\subsection{Nullstellen von Polynomen über $\mathbb{C}$}

\subsection{Nullstellen von Polynomen über $\mathbb{Q}$}

%
% Symmetrische Polynome und die Spur
%
\subsection{Symmetrische Polynome und die Spur}
Der Fundamentalsatz besagt, dass sich ein Polynom in Linearfaktoren
zerlegen lässt.
Multipliziert man die Linearfaktoren aus, entstehen die Koeffizienten
des ursprünglichen Polynoms durch algebraische Ausdrücke aus den
Nullstellen.
In diesem Abschnitt soll untersucht werden, wie diese algebraischen
Ausdrücke entstehen.

%
% Symmetrische Polynome
%
\subsubsection{Symmetrische Polynome}
Die Reihenfolge der Linearfaktoren, in die eine Polynom zerfällt, spielt
keine Rolle.
Jede Reihenfolge führt auf das gleiche ausmultiplizierte Polynom.
Um die Entstehung der Koeffizienten zu verstehen, müssen wir daher nur
algebraische Ausdrücke in den Nullstellen studiert werden, die sich nicht
ändert, wenn man die Nullstellen vertauscht.
Da beim Ausmultiplizieren nur Multiplikation und Addition verwendet
werden, müssen wir also Polynome in mehreren Variablen studieren, die
sich nicht ändern, wenn man die Variablen vertauscht.

\begin{definition}
Ein Polynom $p(X_1,\dots,X_n)$ in den Variablen $X_1,\dots,X_n$
heisst \emph{symmetrisch}, wenn für jede Permutation $\sigma\in S_n$
der Variablen das gleiche Polynome ergibt.
Es gilt also
\[
p(X_{\sigma(1)},\dots,X_{\sigma(n)})
=
p(X_1,\dots,X_n)
\]
für alle Permutationen $\sigma\in S_n$.
\end{definition}

\begin{beispiel}
\label{buch:komplex:fundamental:symmetrisch:bsp:n=2}
Wir bestimmen die symmetrischen Polynome für $n=2$.
Es gibt nur zwei Permutationen für $n=2$, nämlich die identische
Permutation und die Vertauschung $1\leftrightarrow 2$.

Das Polynom $p(X_1,X_2)=X_1$ ist nicht symmetrisch, da die Vertauschung
$1\leftrightarrow 2$ daraus das Polynom $X_2\ne p(X_1,X_2)$ macht.
Erst die Kombination $\sigma_1(X_1,X_2)=X_1+X_2$ ist unter der Vertauschung
invariant.

Symmetrische Polynome zweiten Grades sind zum Beispiel
$\sigma_2(X_1,X_2)=X_1X_2$, aber auch
\begin{align*}
q(X_1,X_2)
&=
X_1^2+X_2^2.
\intertext{Allerdings ist dies kein neues Polynom, denn es ist}
X_1^2+X_2^2
&=
(X_1+X_2)^2 - 2X_1X_2
\\
&=
\sigma_1(X_1,X_2)^2 - 2 \sigma_2(X_1,X_2).
\end{align*}
Das Polynom $X_1^2+X_2^2$ kann also durch $\sigma_1(X_1,X_2)$ und
$\sigma_2(X_1,X_2)$ ausgedrückt werden.

Auch für symmetrische Polynome dritten Grades können wir Beispiele
finden.
In allen Fällen lassen sich diese Polynome zu durch die
$\sigma_k(X_1,X_2)$ ausdrücken:
\begin{align*}
X_1^2X_2 + X_1X_2^2
&=
(X_1+X_2)X_1X_2
=
\sigma_1(X_1,X_2)\cdot\sigma_2(X_1,X_2),
\\
X_1^3 + X_2^3
&=
(X_1+X_2)^3 - 3X_1^2X_2 - 3X_1X_2^2
\\
&=
(X_1+X_2)^3 - 3(X_1^2X_2 + X_1X_2^2)
\\
&=
\sigma_1(X_1,X_2)^3 - 3\sigma_1(X_1,X_2)\cdot \sigma_2(X_1,X_2).
\intertext{Dasselbe gilt für symmetrische Polynome vom Grad 4:}
X_1^2X_2^2
&=
\sigma_1(X_1,X_2)^2
\\
X_1^3X_2+X_1X_2^3
&=
X_1X_2(X_1^2+X_2^2)
\\
&=
\sigma_2(X_1,X_2)\bigl(\sigma_1(X_1,X_2)^2-2\sigma_2(X_1,X_2)\bigr)
\\
&=
\sigma_1(X_1,X_2)^2 \sigma_2(X_1,X_2)
-2
\sigma_2(X_1,X_2)^2
\\
X_1^4+X_2^4
&=
\sigma_1(X_1,X_2)^4
-
4\sigma_1(X_1,X_2)^2 \sigma_2(X_1,X_2)
+8
\sigma_2(X_1,X_2)^2
-
6\sigma_2(X_1,X_2)^2
\\
&=
\sigma_1(X_1,X_2)^4
-
4\sigma_1(X_1,X_2)^2 \sigma_2(X_1,X_2)
+2
\sigma_2(X_1,X_2)^2.
\end{align*}
Keines der symmetrischen Polynome dritten oder vierten Grades enthält
Information, die sich nicht auch schon aus den beiden speziellen
Polynomen $\sigma_1(X_1,X_2)$ und $\sigma_2(X_1,X_2)$ ablesen lässt.
\end{beispiel}

Im Beispiel wurden zwei Polynome gefunden, mit denen sich alle 
anderen symmetrischen Polynome ausdrücken liessen.
Dies sind die sogenannten elemetarsymmetrischen Polynome,
die im nächsten Abschnitt eingeführt werden.

%
% Elementarsymmetrische Polynome
%
\subsubsection{Elementarsymmetrische Polynome}
Wir betrachten das Polynom
\[
p(T,X_1,\dots,X_n)
=
(T+X_1)\cdot\ldots\cdot (T+X_n)
\]
in den Variablen $T,X_1,\dots,X_n$.
Da das Produkt kommutativ ist, ist dieses Polynom symmetrisch in den
Variablen $X_1,\dots,X_n$.
Multipliziert man das Polynom aus und fasst die Terme mit gleicher
Potenz von $T$ zusammen erhält man
\[
\renewcommand{\arraycolsep}{2pt}
\begin{array}{rclclclcl}
p(T,X_1,\dots,X_n)
&=&
T^n
&+&
\sigma_1(X_1,\dots,X_n)
T^{n-1}
&+&
\sigma_2(X_1,\dots,X_n)
T^{n-2}
&+&
\ldots
\\
& &
\ldots
&+&
\sigma_{n-1}(X_1,\dots,X_n) T
&+&
\sigma_n(X_1,\dots,X_n)
\end{array}
\]
mit Polynomen $\sigma_k(X_1,\dots,X_n)$.
Da $p(T,X_1,\dots,X_n)$ symmetrisch ist, gilt
\[
p(T,X_{\sigma(1)},\dots,X_{\sigma(n)})
=
p(T,X_1,\dots,X_n)
\] 
für jede Permutation $\sigma\in S_n$.
Durch Koeffizientenvergleich folgt, dass
\[
\sigma_k(X_{\sigma(1)},\dots,X_{\sigma(n)})
=
\sigma_k(X_1,\dots,X_n)
\]
für jeden Koeffizienten $\sigma_k$, $k=1,\dots,n$, und für jede
Permutation $\sigma\in S_n$.
Insbesondere sind die $\sigma_k(X_1,\dots,X_n)$ alle symmetrische
Polynome.

\begin{definition}[Elementarsymmetrische Polynome]
\label{buch:komplex:fundamental:def:elementarsymmetrisch}
Die Polynome $\sigma_k(X_1,\dots,X_n)\in \mathbb{Z}[X_1,\dots,X_n]$ 
heissen \emph{elementarsymmetrische Polynome}.
\end{definition}

\begin{beispiel}
Für $n=2$ entstehen die elementarsymmetrischen Polynome aus dem
Produkt
\[
(T+X_1)(T+X_2)
=
T^2 +(X_1+X_2)T+X_1X_2
\quad\Rightarrow\quad
\left\{
\begin{aligned}
\sigma_1(X_1,X_2) &= X_1+X_2 \\
\sigma_2(X_1,X_2) &= X_1X_2.
\end{aligned}
\right.
\]
Dies sind genau die Polynome, die im
Beispiel~\ref{buch:komplex:fundamental:symmetrisch:bsp:n=2}
gefunden wurden.
\end{beispiel}

\begin{beispiel}
Für $n=3$ 
kann man aus dem Produkt
\[
(T+X_1)
(T+X_2)
(T+X_3)
=
T^3
+
(X_1+X_2+X_3)T^2
+(X_1X_2+X_1X_3+X_2X_3)T
+
X_1X_2X_3
\]
die elementarsymmetrischen Polynome
\begin{align*}
\sigma_1(X_1,X_2,X_3) &= X_1+X_2+X_3\\
\sigma_2(X_1,X_2,X_3) &= X_1X_2 + X_1X_3 + X_2X_3\\
\sigma_3(X_1,X_2,X_3) &= X_1X_2X_3
\end{align*}
ablesen.
\end{beispiel}

Die Beispiele zeigen auch, dass für beliebiges $n$ die
elementarsymmetrischen Polynome für $k=1$ und 
$k=n$
die besonders einfache Form
\begin{align*}
\sigma_1(X_1,\dots,X_n) &= X_1+\ldots + X_n\\
\sigma_n(X_1,\dots,X_n) &= X_1 \cdots X_n
\end{align*}
haben.
Die anderen elementarsymmetrischen Polynome können 
\[
\sigma_k(X_1,\dots,X_n)
=
\sum_{1\le i_1<i_2<\dots<i_k\le n}
X_{i_1}
X_{i_2}
\cdots
X_{i_k}
\]
geschrieben werden.

%
% Der Fundamentalsatz für symmetrische Polynome
%
\subsubsection{Der Fundamentalsatz für symmetrische Polynome}
Die elementarsymmetrischen Polynome sind in einem gewissen Sinne
die einzigen symmetrischen Polynome, die man braucht.
Dies ist der Inhalt des folgenden Fundamentalsatzes für
symmetrische Polynome: jedes symmetrische Polynom lässt sich
als Polynom in den elementarsymmetrischen Polynomen ausdrücken.

\begin{satz}[Fundamentalsatz für symmetrische Polynome]
Ist $p(X_1,\dots,X_n)\in K[X_1,\dots,X_n]$ ein symmetrisches
Polynom von $n$ Variablen, dann gibt es genau ein Polynom
$q(S_1,\dots,S_n)\in K[S_1,\dots,S_n]$ derart, dass
\[
p(X_1,\dots,X_n)
=
g\bigl(\sigma_1(X_1,\dots,X_n),\dots,\sigma_n(X_1,\dots,X_n)\bigr).
\]
\end{satz}

